\section*{Hoofdstuk \ref{ch:g2:what-animals-eat} --- Wat dieren eten}

\begin{solution}{exo:g2:what-animals-eat:1}
\omterm{def:g2:what-animals-eat:kinds}{Planteneter}; \omterm{def:g2:what-animals-eat:kinds}{vleeseter}; \omterm{def:g2:what-animals-eat:kinds}{alleseter}.
\end{solution}

\begin{solution}{exo:g2:what-animals-eat:2}
(a) \omterm{def:g2:what-animals-eat:kinds}{Planteneter}. (b) \omterm{def:g2:what-animals-eat:kinds}{Vleeseter}. (c) \omterm{def:g2:what-animals-eat:kinds}{Alleseter}. (d) \omterm{def:g2:what-animals-eat:kinds}{Planteneter}.
(e) \omterm{def:g2:what-animals-eat:kinds}{Alleseter}.
\end{solution}

\begin{solution}{exo:g2:what-animals-eat:3}
Platte maalkiezen --- \omterm{def:g1:plants-around-us:plant}{planten} als gras zijn taai en moeten lang
vermalen worden voor het slikken; er is geen prooi om te grijpen.
\end{solution}

\begin{solution}{exo:g2:what-animals-eat:4}
Kort en dik: zaden kraken. Gekromd en gehaakt: vlees scheuren. Lang en
dun: insecten uit bast en spleten pikken.
\end{solution}

\begin{solution}{exo:g2:what-animals-eat:5}
Piepkleine bladluizen --- het lieveheersbeestje eet er enorme aantallen
van, en daarom is het bij tuinders zo geliefd.
\end{solution}

\begin{solution}{exo:g2:what-animals-eat:6}
Gras geeft per hap weinig voeding, dus moet de koe urenlang grazen om
genoeg binnen te krijgen; vlees is rijk, dus kan één goede maaltijd de
vos een hele tijd op de \omterm{def:g1:my-body:parts}{been} houden.
\end{solution}

\begin{solution}{exo:g2:what-animals-eat:7}
Wormen en insecten in de zomer; bessen in de winter, als de grond te
hard is voor wormen. Een breed menu laat een \omterm{def:g2:what-animals-eat:kinds}{alleseter} van voedsel
wisselen als het ene opraakt --- een grote hulp in zware seizoenen.
\end{solution}

\begin{solution}{exo:g2:what-animals-eat:8}
Open antwoord. Meestal: mussen en vinken (korte dikke snavels) pakken
zaden; mezen pikken insecten en zaden; merels prikken naar wormen en
pakken bessen. Het verband tussen snavel en menu uit de propositie
klopt goed.
\end{solution}

\begin{solution}{exo:g2:what-animals-eat:9}
Het menu van een \omterm{def:g2:what-animals-eat:kinds}{vleeseter}: andere \omterm{def:g1:animals-around-us:animal}{dieren}. Stap 1 en 2 --- de tanden
(puntige grijpers) en het jachtgereedschap (klauwen, \omterm{def:g1:the-five-senses:sense}{ogen} naar voren);
stap 4 gaf de naam.
\end{solution}

\begin{solution}{exo:g2:what-animals-eat:10}
De vos eet konijnen, en konijnen eten gras. Geen gras, geen konijnen;
geen konijnen, hongerige vossen. De vos hangt \emph{via} zijn prooi van
het gras af --- maaltijden in de natuur vormen ketens.
\end{solution}
